\section{Durable Trust：Reliability、Security 与透明恢复的复利}

Identity startup 的早期客户必须相信一个小团队能保护组织前门、不会频繁维护停机，并能在事故中承担责任。McKinnon 的 Salesforce 经验提供初始 credibility，但长期信任只能由 production evidence 积累：稳定服务、正确 change、诚实 incident response、customer guidance 和持续平台投入。

\subsection{Trust Flywheel}

可靠服务让客户扩大 deployment；更多应用和用户带来更丰富的 integration/context，也扩大责任；security evidence 与透明 RCA 帮助客户判断 residual risk；客户信任又带来续约、reference 和更深平台采用。飞轮也可能反向运行：一次 poorly handled incident 不只损害当期 SLA，还会让客户减少接入范围，削弱未来 context 和产品价值。

\lecturefigure{15-trust-flywheel.png}{身份平台的采用通过可靠性、安全证据、客户信任与平台深度形成复利。}{本地字幕 27:30--35:20；概念重绘。}

\begin{knowledgebox}{读图：Trust 是证据累积，不是零事故承诺}
Reliable service 证明 front door 可依赖；security evidence 展示 control 和学习能力；customer trust 体现为扩大部署、reference 与 renewal；platform depth 又改善 integration 和 policy context。任何大型系统都无法诚实承诺永不出事故，但可以承诺快速发现、限制 blast radius、公开可行动事实并完成可验证 remediation。
\end{knowledgebox}

\teachervoice{讲者说，早期客户不仅在买一个功能，也在判断“这个人会不会卖完就走、能不能把服务做稳”。他的过往工程经验提供了第一轮信任，但后续每次 outage、breach 和恢复都会重新定价这份信任。}

\subsection{本章小结}

Identity trust 由长期运营证据构成。Reliability、security、transparency 与 customer outcome 互相强化，品牌口号只能放大已有证据，不能替代它。

